\section{Circuit normalization for critical FK maps}
\label{comp:fk-normalization-section}

We now return to the fixed FK parameter $q\in(0,4)$ and use the
quantum volume base $\lambda=4$.  Our deterministic distance denominator
will be a quantile of the length of a circuit surrounding the cone mark.
This circuit test has no prescribed lattice endpoint.  Lower crossing
tests first produce record scales whose earlier quantiles have the
reference power-law upper bound in the limit.  Annulus comparison and
rigidity identify the resulting passage laws.  Their two-sided shell
tests then control the relative cost units, allowing a first-failure
argument to propagate a polynomial ratio bound to all scales.

\subsection{Circuit quantiles and coherent quantum offsets}

Use the circle-average embedding of the Brownian cone, with its mark
at $0$, initially modulo rotations.  Fix three nested buffered winding
bands with strict margins which are concentric round annuli about $0$,
with their coarse-test compact annulus and fixed real-field buffers
contained in $\{1/8<|z|<7/8\}$.  Thus these local tests avoid both
the cone mark and the circle-average pinning circle $|z|=1$.
The following choice
fixes their finite-resolution law once and for all.  Let $P$ be the
bilateral iid raw-word law, and let $\mathsf B$ be the normalized
two-sided Brownian contour law.  For every integer $k\ge0$, choose a
joint law $P_k(dW,dZ^k)$ with marginals $P$ and $\mathsf B$ such that
the resolved contour of $W$, divided by $\sqrt{4^k}$ and viewed at
time unit $4^k$, differs from $Z^k$ by a quantity tending to zero in
probability in the compact-uniform topology.  The invariance principle
and a coupling metric for weak convergence give these choices.  They
are fixed before any quantile or subsequence is selected.

Fix measurable readouts of a Brownian contour for the
intrinsic decorated surface, an arbitrary measurable rotational section
of its circle-average coordinate representative,
its curve, and its field and internal metrics on a countable chart
family.  Only contour-determined readouts are meant here.  Fix also
the deterministic mesh projection convention of
Section~\ref{bg:topology-section}, chosen to commute with rotations,
an ordering of the countably many finite primal circuits, and the open
middle-band winding test.
For example, select vertex representatives by their raw time ordering,
project them by $\eta$, and interpolate a projected edge by a straight
segment; every step of this convention commutes with rotation.
Under $P_k$, project the mesh using the curve readout of $Z^k$.
Define $X_k$ as the least genuine unscaled all-edge length of a
primal circuit whose projected range satisfies this fixed test,
using $\infty$ if the collection is empty.  The first minimizer in
the fixed ordering may be marked.  Countable infima make $X_k$
measurable.  All later appearances of its law refer to this exact
chosen $P_k$ marginal.
The band range and winding requirements are invariant under rotation,
so $X_k$ and its first raw-circuit minimizer do not depend on the chosen
rotational section.  We do not assert that the field in an arbitrary
such section has the standard embedded cone law.

Fix $\alpha\in(0,1)$
and, when $\mathbb P[X_j<\infty]>\alpha$, put
\begin{equation}\label{comp:quantile}
 A_j=\inf\{m\in\mathbb N:m\ge1,\ \mathbb P[X_j\le m]\ge\alpha\}.
\end{equation}
Since the finite costs are integers, this gives
\begin{equation}\label{comp:quantile-slack}
 \mathbb P[X_j\le A_j]\ge\alpha,
 \qquad \mathbb P[X_j<A_j]<\alpha.
\end{equation}
In particular, no continuity of the distribution at its quantile is
required.  The hypothesis \mbox{$\mathbb P[X_j<\infty]>\alpha$} ensures that
this quantile is finite.
Lemma~\ref{comp:fk-coarse-tests} makes the inequality hold for every
sufficiently large $j$; assign arbitrary finite positive values to the
finitely many remaining $A_j$ when defining the normalization sequence.
All quantile-slack arguments below concern the sufficiently large indices.

To compare these quantiles across scales, their defining circuit
variables must occur in one actual word experiment.  The next lemma
preserves each previously fixed marginal while constructing the joint
offset array.  It also retains the field and reference metric in the
same quantum units as the passage observations.

\begin{lemma}[One word and exact offset quantile marginals]
\label{comp:quantile-couplings}
For a sequence of base indices $j\to\infty$, the fixed-offset arrays
can be sampled with one common iid word so that the coordinate at
offset $i$ has exactly the chosen $X_{j-i}$ law.  In every finite
offset family, the auxiliary Brownian contours converge jointly to
the canonical time rescalings of one common contour.  Their actual
joint contour-determined spatial, field and internal-metric readouts
can be retained with this same consistency.
\end{lemma}

\begin{proof}
Sample one word $W\sim P$.  Conditional on that word, draw the entire
countable family $(Z^k)_{k\ge0}$ once from the product of the fixed
kernels $P_k(dZ\mid W)$.  Every offset array reuses the coordinate
$Z^{j-i}$ from this family.  Use the
previously fixed definition of $X_{j-i}$ on that same word and
auxiliary contour.  Its marginal is exactly the one defining the
quantile, even though different offset coordinates are dependent
through $W$.

For $v_i=4^{-i}$, write
$(\mathsf T_i Z)(t)=v_i^{-1/2}Z(v_it)$.  Both $Z^{j-i}$ and
$\mathsf T_i Z^j$ approximate the same resolved word in offset
units.  Thus their compact-uniform difference tends to zero in
probability for every fixed $i$.  A finite union handles any fixed
offset family.  Both arguments of each canonical readout have the
same fixed Brownian marginal.  For a measurable readout into a
Polish space, Lusin's theorem makes it uniformly continuous on a
compact set of marginal probability at least $1-\epsilon$.
Consequently readouts of these two close arguments differ by a
quantity tending to zero in probability, after a loss at most
$2\epsilon$.  Let $\epsilon\downarrow0$, apply the argument to
finite subarrays, and diagonalize over the countable chart family.
This retains the joint field/metric readout, rather than only its
separate marginal laws.  At this stage its arbitrary rotational
sections are identified by quantum-surface equivalence.

The standard residual rotation is now sampled once, uniformly,
independently of the common contour and the entire observation array.
This restores the standard embedded cone law.  The concentric winding
tests, and the Euclidean-advance tests in the concentric compact annulus
used below, are unchanged pointwise; consequently the exact finite
circuit laws and the spatial-advance test laws used in
Lemma~\ref{comp:fk-coarse-tests} are preserved.  An imaginary-field completion
not determined by the contour is likewise adjoined
once, conditional on the entire common limiting contour and this
rotation, independently of the entire actual and reference observation
array under the convention of Proposition~\ref{loc:kernel-closure}.
Its copies at other offsets are transported from that same draw.
More explicitly, let $\phi_i$ map the circle-average representative
at offset $i$ to the common base representative, choosing its orientation
by transport, and put $c_i=\gamma^{-1}\log v_i$.  It is a positive
dilation after this common choice of orientation.  The actual readouts
satisfy
\[
 h_i=h_0\circ\phi_i+Q\log|\phi_i'|-c_i,\qquad
 \eta_i(t)=\phi_i^{-1}\eta_0(v_it),\qquad
 D_i=v_i^{-p}\phi_i^*D_0.
\]
The same identities hold for every internal metric on its transported
domain.  Intrinsic affine covariance of the local kernels then gives
the indexed kernel $\mathcal K_i$ in these units.  These are transported
annuli, which need not coincide as Euclidean sets in the base chart.
In unchanged numerical cost units, if $F_0=CD_0$, the pulled-back
rule satisfies $\phi_i^*F_0=Cv_i^pD_i$.  Thus only the deterministic
quantum-unit factor enters the offset comparison; the random coordinate
dilation supplies no additional normalization.
Unrelated rotations or imaginary-field draws at different offsets
are not used to infer coherence.

The finite variables $Z^k$ serve only to define the quantile and its
projection readouts.  The exact raw-chart resampling identities are
proved before the limiting spatial readout, and no conditional
independence of finite microscopic observations given $Z^k$ is
asserted.  In the limit the contour-determined readouts are the
measurable functions just identified; the additional variables have
the stipulated conditional independence.  These are precisely the
auxiliary-conditioning conditions used by
Proposition~\ref{loc:kernel-closure}.
\end{proof}

Annulus comparison uses local Gaussian field changes.  We next transfer
the indexed passage laws to an outside-pinned Gaussian reference, keeping
their numerical cost units.  The joint offsets then extend the comparison
from ordinary patches of the normalization disk to the whole cone atlas.

\begin{lemma}[Actual Gaussian references and the ordinary cone atlas]
\label{comp:fk-reference-atlas}
The FK local kernels on buffered ordinary patches of the normalized
cone disk transfer to an actual Gaussian field with a fixed pinning
circle outside every observation buffer.  The local field density is
the same for all indexed kernels attached to that field marginal.
If a comparison holds with the same numerical constant for all the
indexed Gaussian laws, it holds on every ordinary cone chart.
Conversely, a strict local improvement on the cone would improve the
Gaussian comparison.  These assertions retain the real additive
constant and the complete joint observation law.
\end{lemma}
\begin{proof}
Inside its normalization disk the cone field is
$G_1-\gamma\log|z|$, where $G_1$ has zero unit-circle average
\cite[Definition~4.10 and the following discussion]{DMS}.
Under the deterministic coordinate map $z\mapsto z/R$, this is
\[
 G_R(z)-\gamma\log|z|+(\gamma-Q)\log R,
 \qquad G_R(z)=G_1(z/R).
\]
On a compact ordinary patch in $B_R\setminus\{0\}$ the displayed
drift has a smooth finite-energy extension supported away from both
$0$ and the pinning circle.  Cameron--Martin gives mutual absolute
continuity with the actual $G_R$ restriction.  To put a prescribed
buffered shell in such a patch, first map it by a fixed positive affine
map into $B_1\setminus\{0\}$.  The pulled-back pinning circle remains
outside the complete buffer; the drift, including $Q\log|\phi'|$,
again has a Cameron--Martin extension.  The pin need not be centered
on the shell.

Any two fixed exterior pins give equivalent restrictions on a smaller
common patch.  Indeed, on an intervening domain $W$ their domain
Markov decompositions have an independent zero-boundary GFF and a
harmonic part.  On $K\Subset W$ each harmonic part admits a
finite-energy cutoff in $W$.  Conditional Cameron--Martin equivalence
with the zero-boundary restriction, followed by integration, proves
the assertion.  No circle average is reset in this comparison.
The independent stationary-phase imaginary field has the same local
law under these affine maps.  Adjoining the same local kernels, with
their genuine order and correct local area marks, preserves each
field density.  If $Z=d\nu/d\mu$, then for any indexed kernel and
event $E$,
\[
 (\nu\otimes\mathcal K_i)(E)
 \le T(\mu\otimes\mathcal K_i)(E)
       +\mathbb E_\mu[Z\mathbf1_{\{Z>T\}}].
\]
The tail is independent of $i$.  Thus uniform probability bounds
transfer by truncating this fixed density; a quantitative exponential
rate under an arbitrary equivalent field law is not asserted.

For the full atlas use the actual common offsets of
Lemma~\ref{comp:quantile-couplings}, including each fixed negative
offset.  Put $c_v=\gamma^{-1}\log v$ and
\[
 r_v=\sup\{r>0:h_r(0)+Q\log r=c_v\},\qquad
 h^{(v)}=h(r_v\,\cdot)+Q\log r_v-c_v.
\]
The circle-average cone law and its quantum scaling
\cite[Proposition~4.13(i)]{DMS} give the original joint background law
at every fixed $v$.  The imaginary field is pulled back from the same
sample; its stationary phase and independence of the real field
preserve this joint law.  The matching raw time, height, cost, and
reference arrays were transformed together before extraction.  Their
intrinsic readouts therefore obey
\[
 F^{(v)}_{U}=v^{-p}\phi_v^*F_{\phi_v(U)},\qquad
 D_{h^{(v)}}(\cdot,\cdot;U)
       =v^{-p}\phi_v^*D_h(\cdot,\cdot;\phi_v(U)),
 \quad \phi_v(z)=r_vz,
\]
and the local kernel in the new units is $\mathcal K^{(v)}$.
These are identities for the transported observations, with fresh
reference kernels identified by Proposition~\ref{loc:kernel-closure}.
They are not field-shift identities for an individual passage sample.

Set $H(r)=h_r(0)+Q\log r$.  The radial cone construction gives
$H(r)\to-\infty$ as $r\downarrow0$, and $H$ is continuous.
Hence $\sup_{0<r\le R}H(r)<\infty$ for each finite $R$.
Since $c_{4^m}\to\infty$, it follows that $r_{4^m}\to\infty$.
For each fixed $m$, the comparisons under the indexed Gaussian kernel
transfer to every ordinary patch inside the normalized cone disk.
Pulling them back cancels the two factors $4^{-mp}$ and covers
$B_{r_{4^m}}\setminus\{0\}$.  A countable intersection over $m$
and a countable chart exhaustion gives the whole ordinary cone atlas.
This uses an exact joint quantum-unit identity, without conditioning
on a random dilation and then applying local absolute continuity.
For the reverse assertion, transfer a local comparison on ordinary
patches of the unit disk to the Gaussian reference by mutual absolute
continuity and intrinsic affine transport.  Lemma~\ref{comp:comparison-transport}
then preserves its numerical constant at every fixed quantum index.
\end{proof}

\begin{lemma}[Noncollapse of a quantile-normalized upper rule]
\label{comp:quantile-noncollapse}
Suppose a sequence normalized by $A_j$ has a certificate limit, winding
realizations pass from the narrow band to the middle band with arbitrary
cost loss, and the limiting rule can follow every rectifiable path at
zero cost with open ports.  If the reference narrow band has a finite
rectifiable surrounding circuit almost surely, these assumptions
contradict Equation~\eqref{comp:quantile-slack}.
\end{lemma}

\begin{proof}
Cover a reference surrounding circuit by finitely many transverse port
tests in the narrow band.  Zero upper comparison gives certificates of
arbitrarily small total cost; robust transverse intersections stitch
them into a primal winding circuit.  Exhaust its finite number of
ports, reference length, and strict margins.  For every $\delta>0$, this
gives, with limiting probability at least $1-\delta$, a middle-band
circuit of cost less than $A_j/2$ for all sufficiently large $j$.
Choose $\delta<1-\alpha$.  Its probability is then strictly greater
than $\alpha$, contradicting
$\mathbb P[X_j<A_j]<\alpha$.
\end{proof}

The preceding lemma excludes a zero upper comparison after an upper
comparison has been constructed.  A single quantile does not by itself
give tight lower band costs, a finite upper comparison, or any estimate
for nearby discrete endpoints.

\begin{lemma}[Concrete finite-resolution tests]
\label{comp:fk-coarse-tests}
Assume the framework's contour coupling, marked certificate extraction,
and robust primal projection topology on compact sets.  The winding
variables $X_j$ can then be chosen so that
$\mathbb P[X_j<\infty]\to1$ and $A_j\to\infty$.
There is a fixed finite collection of compact spatial-advance tests
that satisfies the crossing-test hypotheses of
Proposition~\ref{comp:initial-growth}.
\end{lemma}

\begin{proof}
Keep the three concentric winding bands fixed above, and choose a
concentric compact annulus $K$, bounded away from the cone mark,
containing the widest band in its interior.  Finite disk covers of a
slightly larger annulus supply all required local chart patches.  Choose
$\delta>0$ smaller than the Euclidean diameter forced by winding
around the inner hole.  A deterministic polygonal loop in the
narrowest band has an approximate primal lift by the projection
topology.  Its lift is a finite graph walk, stays in the middle band,
and has nonzero winding.  Extract a simple surrounding subcircuit.
Thus, on the events where the projection error is smaller than the
fixed margins, $X_j$ is finite.  These events have probability tending
to one.  Time windows are first enlarged to contain the preimages of
the annulus; transience gives this with arbitrarily high probability.
No cost bound for the lift is asserted.

Define $Y_j$ to be the infimal number of edges of a primal path whose
projected range lies in $K$ and whose projected endpoints are at
Euclidean distance at least $\delta$.  Use $\infty$ for an empty
collection.  Slightly larger compact confinements and slightly smaller
separations provide strict versions.  The contour coupling and the
closed chord relation make the largest projected diameter of a single
edge tend to zero, uniformly on the relevant compact time windows.
A path of at most $k$ edges has diameter at most $k$ times this
quantity.  Therefore, for each fixed $k$,
$\mathbb P[Y_j\le k]\to0$.  Every winding circuit tested by $X_j$
contains an admissible path for $Y_j$, so $Y_j\le X_j$ and
$A_j\to\infty$.

At any proposed cost unit $b_j$, failure of the test
$\mathbb P[Y_j\ge b_j]\ge1-\eta$ retains, with probability at
least $\eta$, a path of normalized cost at most one and endpoints
separated by at least $\delta$.  The endpoints and range lie in fixed
compact sets.  Marked extraction retains a certificate with that
separation and finite cost.  First choose a sufficiently large time
window that its lost probability is less than $\eta/2$, then pass to
the extraction and exhaust windows.  The fixed spatial margins retain
the separated endpoints and compact confinement.

Choose a buffered annulus inside the interior of $K$ with radial
crossing separation greater than $\delta$.  Its primal traversal
tests are included among these spatial-advance tests.  Translations
and the fixed affine changes needed to place the annulus use a finite
collection of slightly enlarged versions of $K$.  The common local
kernels transfer their high-probability tests to the fixed GFF
reference: the field change of density is fixed, so
Equation~\eqref{comp:ann-reference-change} allows $\eta$ to be chosen
small enough for the annulus threshold.  Uniform bounded-edge advance
gives the required initial runs, and the retained finite-cost witness
gives the required first-failure contradiction.  These are precisely
the test conditions of Proposition~\ref{comp:initial-growth}.
\end{proof}



\subsection{Deterministic scale arguments}

We initialize a positive lower barrier, obtain rigidity at record scales,
control relative quantiles, and propagate the comparison by excluding a
first failure.

\begin{proposition}[Lower-barrier initialization]
\label{comp:initial-growth}
Assume the following assertions for arbitrary deterministic cost
normalizations and the same fixed buffered tests.
\begin{enumerate}[label=(\alph*)]
\item Every sequence of finite-cost crossing witnesses has a
subsequential certificate retaining a fixed positive spatial advance.
The common-kernel and annulus transfer properties hold in these limits.
\item At unscaled cost, the largest projected advance of every path
using at most $k$ edges tends to zero in probability, for each fixed
$k$, uniformly in the compact testing region.
\item A winding circuit used in $X_j$ contains a subpath tested by a
fixed finite collection of these crossing tests.  The probability and
cost thresholds have strict margins, so failure of their high-probability
lower test retains a finite-cost witness with probability bounded away
from zero.
\end{enumerate}
Then, for the finite circuit quantiles in Equation~\eqref{comp:quantile},
\[
 \limsup_{j\to\infty}j^{-1}\log_4A_j\ge p.
\]
\end{proposition}

\begin{proof}
Fix $0<\varepsilon<p$.  Start at $j_0$ with denominator one and use
$b_j=4^{(p-\varepsilon)(j-j_0)}$ thereafter.  Choose the finite lower
test's failure probability smaller than the annulus threshold,
$\alpha$, and $1-\alpha$, distributing it among the finitely many crossing
geometries.  By (b), every fixed initial run passes as $j_0\to\infty$.

Suppose there are arbitrarily late starts for which a first failure
$f(j_0)$ occurs.  Then $f(j_0)-j_0\to\nobreak\infty$.  Extract at the failed
scale in units $b_{f(j_0)}$.  At each fixed preceding offset $i$ the
lower crossing threshold, expressed in reference metric units, is
\[
 \frac{b_{f(j_0)-i}}{b_{f(j_0)}4^{-ip}}
 =4^{\varepsilon i}.
\]
For any $T<\infty$, all sufficiently large labels therefore have
barriers of at least $Tw_i$, with reference surrounding costs at most
$A_0w_i$ on a sufficiently high probability regularity event.  Annulus
transfer and Corollary~\ref{comp:annulus-comparison} force every
nonconstant retained certificate to have cost at least $(T/A_0)D$
between its endpoints.  Take the countable intersection over integer
$T$.  This excludes every finite-cost certificate with positive spatial
advance, contrary to the failed-test witness from (a) and (c).

Consequently a sufficiently late start passes forever.  With the chosen
probability margins, (c) gives $\mathbb P[X_j<b_j]<\alpha$ at every
subsequent level, and hence $A_j\ge b_j$ (integer rounding changes this
by at most one).  This yields the stated lower growth after letting
$\varepsilon\downarrow0$.
\end{proof}

For the FK maps, Lemma~\ref{comp:fk-coarse-tests} supplies these strict
crossing tests and their retained finite-cost failure witnesses.

\begin{lemma}[Record scales]\label{comp:records}
If $A_j\in(0,\infty)$ and
$\limsup_{j\to\infty}j^{-1}\log_4 A_j\ge p$, then there are
$r_k\uparrow\infty$ and $t_k\uparrow p$, $t_k<p$, such that
\[
 \frac{A_{r_k-i}}{A_{r_k}}\le 4^{-t_k i},\qquad 0\le i\le r_k.
\]
\end{lemma}

\begin{proof}
For every fixed $t<p$, the sequence $4^{-tj}A_j$ is unbounded and has
arbitrarily late record maxima.  At a record $r$,
$4^{-t(r-i)}A_{r-i}\le4^{-tr}A_r$.  Choose successively larger records
for a sequence $t_k\uparrow p$.
\end{proof}

\begin{lemma}[Continuum comparison at the cone mark]
\label{comp:fk-marked-extension}
Suppose the FK family has a common finite upper comparison at all
fixed quantum indices on the ordinary cone atlas, and arbitrarily
small realizable primal separating circuits at every ordinary point,
with costs tending to zero and strict winding margins.  Then almost surely
the cone mark is surrounded by realizable primal circuits whose
diameters and normalized costs tend to zero.  The port comparison,
continuum concatenation, and any intrinsic lower comparison extend
across that mark with the same constants.
\end{lemma}
\begin{proof}
In the fixed narrow reference annulus choose an almost surely finite
rectifiable surrounding circuit, with a finite length bound $Y$.
The upper comparison on this ordinary annulus, finite transverse
port tests, and small primal separators at its ordinary points
realize a winding primal circuit in the middle annulus, with cost
at most $CY+1$.  Only this fixed annular construction is used here.

Apply it in each of the exact units $v=4^{-m}$ retained in
Lemma~\ref{comp:quantile-couplings}.  In base units the circuit has
cost at most $v^p(CY_v+1)$ and lies in a fixed multiple of
$B_{r_v}(0)$, where $Y_v$ has the same finite marginal law as $Y$.
The radial construction in \cite[Definitions~4.5 and~4.10,
Remark~4.4]{DMS} gives $H(r)\to+\infty$ as $r\to\infty$ and
$H(r)\to-\infty$ as $r\downarrow0$, since $Q-\gamma>0$.
Consequently $\inf_{r\ge\varepsilon}H(r)>-\infty$ for every
$\varepsilon>0$, and $r_v\to0$ as $v\downarrow0$.
Select a deterministic subsequence $m_k$ so that the probabilities
of a radius larger than $2^{-k}$ and of a cost larger than $2^{-k}$
are each summable in $k$.
Borel--Cantelli gives the asserted circuits, with strict winding
margins and simultaneous realizations.  No independence of the
different offsets is required.

For a compact reference path meeting $0$, use the first and last
crossings of one of these circuits.  The exterior initial and final
subpaths are compactly away from $0$, so their port upper bounds
are already known.  Join their realizations along the primal circuit,
whose extra cost tends to zero.  If an endpoint is $0$, choose the
circuit inside its prescribed open endpoint port.  This proves the
same upper port comparison for all paths.  The circuits at $0$ and
the existing circuits at ordinary points now give
Lemma~\ref{comp:port-splicing} on the whole chart.

For the lower comparison, retain the ordered prefix and suffix
outside $B_\varepsilon(0)$ of a realizing certificate.  The omitted
middle cost is nonnegative.  The known lower bounds for the two
exterior pieces give the lower bound for its endpoints up to
$c\operatorname{diam}_D(B_\varepsilon(0))$, which tends to zero
by Lemma~\ref{common:cone-mark-continuity}.  The same argument with just one exterior
piece treats an endpoint at $0$; if both endpoints are $0$ the
bound is zero.  First extract these markings with open margins,
then let $\varepsilon\downarrow0$.  This proves the ambient lower
bound everywhere, and Lemma~\ref{comp:intrinsic-lower} makes it
intrinsic.  This is a statement about continuum certificates.
\end{proof}

\begin{proposition}[Comparison at exact-power records]
\label{comp:exact-power-comparison}
Assume the certificate, spatial-kernel, coordinate-transport, and
extraction properties used above.  Choose the quantile level $\alpha$
sufficiently close to one for Lemma~\ref{comp:ann-coverage}, including
reference crossing lower bounds in its local event.  Suppose a joint
extraction in units $A_j$ has
\[
 \limsup_j A_{j-i}/A_j\le K4^{-ip},\qquad i\ge0,
\]
where $K<\infty$ is deterministic.  Then its port rule has a finite
positive deterministic optimal upper comparison $C$.  The same numerical
upper bound is valid at every fixed quantum index and on the full cone
atlas before changing its own cost normalization.  If the family
satisfies Definition~\ref{common:admissible},
then this rule satisfies $F_U=CD_U$.
\end{proposition}

\begin{proof}
At label $i$ the wide-band upper circuit test has, with probability
at least $\alpha$ and arbitrary cost slack, a surrounding circuit
of cost at most $K4^{-ip}$.  Include a reference lower crossing bound
$a_0 4^{-ip}$, with failure probability chosen before $\alpha$.
The common coordinate kernels make this the corresponding fixed-index
test under the local reference field law; fixed changes of density
are handled by Equation~\eqref{comp:ann-reference-change}.
Lemma~\ref{comp:ann-coverage} supplies arbitrarily fine circuits at
every point of each compact set.  Apply the actual-primal-circuit
version of Lemma~\ref{comp:circle-chain}, charging a reference path's
complete traversals.  It gives upper factor $K/a_0$, up to the fixed
geometric constants.  The circuit costs also tend uniformly to zero,
so Lemma~\ref{comp:port-splicing} applies.

The same bound holds for every fixed forward index $v_0=4^m$.
Indeed, its sufficiently fine indexed laws satisfy
\[
 (\mathcal K^{(4^m)})^{(4^{-\ell})}
       =\mathcal K^{(4^{m-\ell})},\qquad \ell\ge m.
\]
This tail belongs to the backward family with the same shell-test
bounds.  Deleting finitely many labels does not change the favorable
density in arbitrarily late windows or the resulting comparison
ratio.  Thus one finite bound works for every fixed $m$, without a
rate uniform in $m$.  The same fine annuli provide the small primal
circuits at ordinary points at every fixed index.
Lemma~\ref{comp:fk-reference-atlas} transfers these assertions
to the ordinary cone atlas, and Lemma~\ref{comp:fk-marked-extension}
extends it across the mark.  The upper optimum cannot be zero by
Lemma~\ref{comp:quantile-noncollapse}.  Its optimal constants have the
same numerical values in the indexed laws by the reverse transfer in
Lemma~\ref{comp:fk-reference-atlas}.  Theorem~\ref{common:identification}
gives equality on ordinary charts, and the same marked extension
gives the stated intrinsic conclusion.
\end{proof}


The preceding comparison identifies a limiting passage law.  We now
return to the winding tests that chose its distance unit: their quantile
slack bounds both the multiplier and the quantiles at every fixed offset.

\begin{lemma}[Quantile bounds in a rigid law]\label{comp:rigid-quantiles}
Suppose a joint extraction, including any fixed finite set of quantum
offsets, has $F_U=CD_U$ in base cost units, where $C\in(0,\infty)$.
Assume the narrow, middle, and wide winding tests used in
Equation~\eqref{comp:quantile} have the following strict passage rules:
upper realizations in the narrow band pass to the middle band, and
middle-band circuits have retained certificates in the wide band.
Use the fixed reference law after each quantum-unit change.
Then there are $0<m_*<M_*<\infty$, independent of this rigid law and
the fixed offset $i$, such that
\begin{equation}\label{comp:rigid-ratio-bounds}
 m_*C4^{-ip}\le
 \liminf\frac{A_{j-i}}{A_j}
 \le\limsup\frac{A_{j-i}}{A_j}
 \le M_*C4^{-ip}.
\end{equation}
In particular $M_*^{-1}\le C\le m_*^{-1}$ when the base is normalized
by its own quantile.
\end{lemma}

\begin{proof}
In the wide reference band, every winding circuit has two points at
a fixed positive Euclidean separation.  By $D$-continuity on the
compact band, the infimum $Y_-$ of the distances of all such pairs is
strictly positive almost surely.  In the narrow reference band there
is an almost surely finite reference rectifiable surrounding circuit;
let $Y_+<\infty$ be any measurably chosen upper length bound.  Choose
$m_*>0$ sufficiently small and $M_*<\infty$ sufficiently large that
\[
 \mathbb P[Y_->2m_*]>1-\alpha/2,\qquad
 \mathbb P[Y_+<M_*/2]>(1+\alpha)/2.
\]
The reference laws of these variables are the same at every fixed
offset in their own quantum units.

If $A_{j-i}/A_j$ had a subsequential limit below
$m_*C4^{-ip}$, the probability-$\alpha$ quantile upper event would
retain a wide-band certificate below this cost with probability at
least $\alpha$.  The lower comparison and the $Y_-$ bound permit
this with probability less than $\alpha/2$, a contradiction.  This
argument also rules out a zero limiting ratio.

Conversely, on the event $Y_+<M_*/2$, upper comparison and transverse
port stitching realize a middle-band circuit below
$M_*C4^{-ip}A_j$, with strict cost and probability margins.
The quantile definition forces $A_{j-i}$ below this bound for all
sufficiently large $j$ along the extraction.  It also rules out an
infinite limiting ratio.  All costs and probabilities were separated
from their thresholds, so no continuity at a quantile atom was used.
Set $i=0$ to obtain the bounds on $C$.
\end{proof}

\begin{lemma}[Polynomial bound after propagation]
\label{comp:polynomial-algebra}
Suppose there are $J,m\in\mathbb N$, $R\ge1$, and $s>0$ such that
for $j\ge J+m$,
\[
 A_{j-1}/A_j\le R,\qquad A_{j-m}/A_j\le4^{-sm}.
\]
Then, after enlarging a finite deterministic constant $K$,
\begin{equation}\label{comp:polynomial-normalization}
 \frac{B(n')}{B(n)}\le K(n'/n)^s\quad(1\le n'\le n),
 \qquad B(n)=A_{\lfloor\log_4 n\rfloor}.
\end{equation}
In particular $B(n)\to\infty$.
\end{lemma}

\begin{proof}
For indices beyond the finite exceptional range write $j-q=am+b$ with
$0\le b<m$.  Iteration gives
$A_q/A_j\le4^{-sam}R^b\le(R4^s)^{m-1}4^{-s(j-q)}$.
The finitely many earlier indices can be joined through $J+m$ and
absorbed in the constant.  If $q=\lfloor\log_4 n'\rfloor$ and
$j=\lfloor\log_4 n\rfloor$, then
$4^{q-j}\le4n'/n$.  This proves Equation~\eqref{comp:polynomial-normalization}.
Taking $n'=1$ gives $B(n)\ge K^{-1}B(1)n^s$.
\end{proof}

\begin{proposition}[First-failure propagation]
\label{comp:propagation}
Assume the preceding framework and common-kernel hypotheses, including
closure under varying-law extraction with all fixed-offset reference
experiments retained as in Proposition~\ref{loc:kernel-closure}.
Assume that the robust lower tests of
Proposition~\ref{comp:initial-growth} and the finite catalogue of
Proposition~\ref{comp:rigidity-implication} have been verified for the
FK certificate family.  Assume also that the strict winding transfer
rules of Lemma~\ref{comp:rigid-quantiles} hold at all fixed offsets.
Then there are deterministic $J,m,R,K$ and $s=p/2>0$ such that all
sufficiently large scales satisfy adjacent ratio bounds in $[R^{-1},R]$,
$A_{j-m}/A_j\le4^{-sm}$, and the high-probability lower shell tests.
Consequently Equation~\eqref{comp:polynomial-normalization} holds.
\end{proposition}

\begin{proof}
\smallskip\noindent\emph{A neighborhood of the rigid laws.}
We describe the required law neighborhood, since a first-failure
argument needs more than a pointwise quantile bound.  Use the product
topology of the full marked certificate records, the contours and
the reference internal metrics on countably many buffered domains.
Include the auxiliary reference experiments that ensure closure of
the common-kernel identities by Proposition~\ref{loc:kernel-closure}.
The cost coordinates are compactified;
the reference marginal is fixed and tight.  Thus these joint laws are
precompact under the assumed framework.

Let $\mathcal R$ be the closure of the quantile-normalized rigid laws
arising in these extractions.  Lemma~\ref{comp:rigid-quantiles} puts
their deterministic factors in a fixed compact subset of $(0,\infty)$.
The closure still consists of rigid laws.  Indeed, along a convergent
subsequence their deterministic factors converge.  Marked extraction
retains the intrinsic lower comparison.  Finite port realizations and
the uniform reference moduli retain the upper comparison and its
small-circuit splicing.  These arguments are the same two directions
used in Proposition~\ref{comp:equality-tangent}.  Hence $\mathcal R$
is compact and its laws all have the same uniform comparison bounds.

Choose coarse shell constants $0<m_0<M_0<\infty$ so loosely that every
law of $\mathcal R$ passes the tests in
Proposition~\ref{comp:relative-units} with strict cost, probability,
and confinement margins.  Such constants are obtained from the fixed
reference law and the compact interval of factors.  Upper tests ask
for finitely many strict port or winding realizations in an open
margin; lower tests exclude a closed set of bounded-cost certificates
with separated endpoints in compact confinement.  Retaining a
slightly wider band and using a slightly weaker cost threshold makes
these tests stable under nearby laws.  More precisely, after the
geometric slack has been fixed their good events are open in the
retained data topology.  For an upper test this is the existence of
a certificate satisfying strict cost, port, and open-confinement
conditions.  An upper winding test uses a finite marked recipe of
robust transverse port crossings; it does not ask for equality of
two continuously perturbed loop endpoints.  For a lower test the
good event is disjointness from a closed
bounded-cost certificate set with compact confinement and separated
endpoints.  For each $\mu\in\mathcal R$, choose a bounded continuous
function $f$ below the indicator of each such good event, with
$\int f\,d\mu>1-\epsilon/2$.  The neighborhood
$\{\nu:\int f\,d\nu>1-3\epsilon/4\}$ has closure contained in
the laws assigning that good event probability at least
$1-3\epsilon/4>1-\epsilon$.  Finite time and chart cutoffs lose
arbitrarily little probability, since their continuum marginal
is fixed.  Intersect the finitely many neighborhoods for the tests.
Compactness then gives a finite union $\mathcal O$ covering
$\mathcal R$, whose closure passes all the weaker coarse tests.

By Lemma~\ref{comp:rigid-quantiles}, every rigid joint limit has all
fixed ratios in
\[
 \frac{m_*}{M_*}4^{-ip}
 \le A_{j-i}/A_j\le
 \frac{M_*}{m_*}4^{-ip}.
\]
Choose $R$ strictly larger than both adjacent ratio bounds and their
reciprocals.  Then choose $m$ large enough that
$(M_*/m_*)4^{-mp}<\tfrac12 4^{-ms}$, with $s=p/2$.
Call a finite scale successful if its own-unit law belongs to
$\mathcal O$ and its adjacent and $m$-step ratio inequalities hold
with these thresholds.

\smallskip\noindent\emph{Record scales start long successful runs.}
Initialization and Lemma~\ref{comp:records} supply record indices
$r_k\to\infty$ with complete preceding tails
$A_{r_k-i}/A_{r_k}\le4^{-t_ki}$, where $t_k\uparrow p$ and
$t_k\ge s$.  Proposition~\ref{comp:exact-power-comparison} and
rigidity apply to every joint record extraction.  Apply the same
argument on a fixed finite forward enlargement, initially retaining
the record's cost units.  At forward offset $a$, its fine past
offsets satisfy
$A_{r_k+a-\ell}/A_{r_k}\le4^{-t_k(\ell-a)}$ for fixed
$\ell\ge a$, whose limit is $4^{pa}4^{-p\ell}$.
These bounds first give the enlarged comparison in the old record
units.  A zero upper factor would contradict the retained old
$X_{r_k}$ quantile by Lemma~\ref{comp:quantile-noncollapse}, before
any bound on the new quantile is needed.
Lemma~\ref{comp:rigid-quantiles} makes its
own quantile ratios positive and finite, so changing to those units
gives laws in $\mathcal R$.  All numerical inequalities pass with
strict margins.  Therefore every fixed finite run after a sufficiently
late record is successful.  Diagonalization gives successful initial
runs of lengths tending to infinity.

\smallskip\noindent\emph{The successful tail and the first-failure extraction.}
Before a first failure, the full record tail remains useful.  If all
levels from $r$ to $j$ are successful and $r$ has its record tail,
iteration of the $m$-step inequalities and the remaining at most
$m-1$ adjacent inequalities gives
\[
 \frac{A_q}{A_j}\le K_0 4^{-s(j-q)},\qquad 0\le q\le j,
 \qquad K_0=(R4^s)^{m-1}.
\]
For $q<r$, combine the same bound at $r$ with the record tail.

Suppose unsuccessful levels occur arbitrarily late.  Let $f_k$ be
the first failure after $r_k$, and set $j_k=f_k-1$.  The initial
success runs imply $j_k-r_k\to\infty$.  Extract in units $A_{j_k}$.
Every fixed preceding own-unit law lies in $\overline{\mathcal O}$,
and the adjacent bounds ensure that each fixed relative quantile
has a positive finite limit.  Write
$b_i=4^{ip}\lim_k A_{j_k-i}/A_{j_k}$, and let $G_i$ denote the
limiting passage rule at offset $i$ in its own cost units.  For fixed
$i,\ell\ge0$, viewed from base $i$, the cost unit at offset $i+\ell$ contributes the ratio
$A_{j_k-i-\ell}/A_{j_k-i}$.  Dividing by the reference length unit
$4^{-\ell p}$ gives
\[
 4^{\ell p}\lim_k\frac{A_{j_k-i-\ell}}{A_{j_k-i}}
       =\frac{b_{i+\ell}}{b_i}.
\]
The limiting denominators $b_i$ are positive by the preceding adjacent
bounds.
The common-word coupling of Lemma~\ref{comp:quantile-couplings}
therefore makes $(b_{i+\ell}/b_i)G_{i+\ell}$ the actual physical
offset family of each base $i$.  The assumed common-kernel framework
and its annulus transfer,
Lemma~\ref{comp:ann-coverage}, supply the physical compatibility
hypothesis of Proposition~\ref{comp:relative-units}.
Proposition~\ref{comp:relative-units} makes all $b_i$ uniformly
comparable.  Consequently the past ratios satisfy an exact-power
upper bound.  Proposition~\ref{comp:exact-power-comparison} and
rigidity make the base law rigid.

\smallskip\noindent\emph{The failed scale in preceding units.}
Examine the failed level in the preceding units before normalizing it.
It is only a fixed time enlargement.  Its cost records are still
available, and the fine past offsets just used give a finite upper
comparison there.  This upper optimum is positive: a zero upper
comparison would restrict to the retained old winding annulus and
contradict Lemma~\ref{comp:quantile-noncollapse} for the old
$X_{j_k}$ quantile.  Rigidity is therefore available before any division
by the failed scale's own normalization.
Lemma~\ref{comp:rigid-quantiles} then makes
$A_{f_k}/A_{j_k}$ positive and finite in the limit.  The failed
own-unit law therefore belongs to $\mathcal R$, and all its ratio
tests lie strictly within the successful bounds.  Since $\mathcal O$
is open, it is successful for large $k$, a contradiction.
All large scales are successful.  Finally
Lemma~\ref{comp:polynomial-algebra} gives the stated polynomial bound.
\end{proof}

\begin{proposition}[Verification for the FK certificate family]
\label{comp:fk-application}
The bilateral FK certificate family constructed in
the preceding framework and localization sections
satisfies the hypotheses of Proposition~\ref{comp:propagation}.
\end{proposition}

\begin{proof}
\emph{Initial tests and local kernels.}
Lemma~\ref{comp:quantile-couplings} fixes the exact quantile marginals
and their joint offset readouts before any normalization extraction.
It also verifies the contour-determined and independent-augmentation
conventions required by Proposition~\ref{loc:kernel-closure}.
Lemma~\ref{bg:certificates} supplies full marked records, realizability,
and additive cost subdivisions.  Proposition~\ref{bg:projection-topology}
supplies primal winding and intersection, the small projected advance
of a bounded number of edges, and the approximate mesh lifts.
Lemma~\ref{comp:fk-coarse-tests} therefore gives finite circuit
quantiles, their divergence, and the concrete separated-endpoint
lower tests.  Its nested compact confinements and the three winding
bands give the strict transfer rules at every fixed time offset.

Theorem~\ref{loc:spatial-kernels} supplies restriction-compatible
spatial kernels with arbitrarily small real-field buffers, a possibly
larger imaginary buffer, and finite genuine traversal order.
The fixed-chart reference experiments and
Lemma~\ref{loc:countable-selection} are retained jointly for all
countable chart refinements and fixed quantum offsets.
Proposition~\ref{loc:kernel-closure} proves closure for the resulting
admissible laws, including the further varying-law and changing-unit
extractions used in the equality tangent and first-failure argument.
Its diagonal discrete approximation retains every prescribed finite
array and sends every fixed rescaled time unit to infinity in lattice
units.  The exact density calculation and local limit theorem apply
anew to that diagonal array and identify its kernels with its own
reference kernels.  Its retained actual joint field/internal-metric
marginal is unchanged.  This verifies the closure hypothesis without
inferring conditional independence from weak convergence alone.
Intrinsic coordinate
transport is part of Theorem~\ref{loc:spatial-kernels};
Equations~\eqref{comp:ann-index-definition} and
\eqref{comp:ann-two-units} specify all changes of cost and volume
units.  Corollary~\ref{loc:field-tilt} supplies their density transfer
and the correct resampling of local area marks.  The independent
real and imaginary GFF inputs and their local equivalences come from
Theorem~\ref{bg:continuum-inputs}.

In particular, the indexed family in
Equation~\eqref{comp:ann-index-definition} uses actual discrete
offset experiments at times $Nv_i$ and cost denominators
$B_Nv_i^p$, retained jointly before applying the fixed local
field-density changes.  This checks admissibility of the indexed
laws required by Proposition~\ref{loc:kernel-closure}; it is not
deduced solely from the algebraic kernel formula.

\smallskip\noindent\emph{Annulus comparison.}
Consequently Lemma~\ref{comp:ann-coverage} applies.  It proves the
annulus transfer needed for initialization, exact-power comparison,
and amplification.  Small circuits give port splicing before any
endpoint continuity of $F_U$ is used.  The exact tail reindexing in
Proposition~\ref{comp:exact-power-comparison} gives one numerical
upper bound at every fixed quantum index.  The actual atlas in
Lemma~\ref{comp:fk-reference-atlas}, followed by
Lemma~\ref{comp:comparison-transport}, then identifies the optimal
constants of the whole joint law with those of its indexed members.  The metric and confluence
inputs for equality tangents are exactly
Theorem~\ref{bg:metric-inputs}.

\smallskip\noindent\emph{The common interface and forcing.}
We have verified Definition~\ref{common:admissible}(i)--(iv) using
marked compact certificate records, the spatial product kernels,
actual indexed experiments, and the diagonal closure theorem.
For the FK records the compactified cost and marked path coordinates
provide the required tightness before comparison; after comparison,
reference metric moduli and small-circuit joining give the uniform
continuum moduli used in its further extractions.  The edge-cost
condition follows from the divergence of the quantile units.  The
finite genuine order satisfies (vi) by
Lemma~\ref{comp:local-order-certificate}.  Once the annulus argument
has provided the finite positive upper comparison, small circuits
and Lemma~\ref{comp:port-splicing} supply (v), before any theorem
about a prescribed discrete endpoint is used.

Lemma~\ref{comp:fk-reference-atlas} supplies the actual outside-pinned
Gaussian law and the countable cone atlas.  In particular, all recipes
in the forcing argument are Cameron--Martin directions for one fixed
Gaussian law, and the same local kernels are used with correctly
resampled area marks.  The reference inputs are
Theorem~\ref{bg:metric-inputs}.  Therefore
Lemma~\ref{common:forcing-verification} supplies the complete catalogue,
its deterministic gate budgets, unselected completion followed by
likelihood clipping, source weight and observable output count.  Its disk-profile bound is
Lemma~\ref{common:disk-profiles}; it retains the real preservation
patch inside the denominator collar and allows a larger unchanged
imaginary buffer.  Theorem~\ref{common:identification} applies to
every equality-tangent family, and
Lemma~\ref{comp:fk-marked-extension} gives the conclusion across the
cone mark.

\smallskip\noindent\emph{Rigidity and propagation.}
Rigidity is therefore available for every family used at records
or at a proposed first failure.  The remaining inputs of
Proposition~\ref{comp:propagation} were verified above, so its
first-failure argument applies to the actual FK normalizations.
\end{proof}

\begin{corollary}[Bilateral comparison]
\label{comp:bilateral-comparison}
For the fixed FK parameter, every sequence of bilateral time units
tending to infinity has a further
joint extraction, normalized by $B(n)$, with
\[
 F_U=C D_h(\cdot,\cdot;U)
\]
simultaneously in all strict intrinsic charts.  The multiplier $C$
is deterministic for that extraction and belongs to a fixed compact
subset of $(0,\infty)$.  Upper comparisons follow paths inside any
open clearance; lower comparisons apply to all compactly confined
certificates and use the intrinsic metric of their confinement.
Also $B(n)\to\infty$ and
Equation~\eqref{comp:polynomial-normalization} holds.  This conclusion
concerns crossing data and infimal continuum costs, and supplies no
uniform estimate for arbitrary discrete endpoints.
\end{corollary}

\begin{proof}
Proposition~\ref{comp:fk-application} permits application of
Proposition~\ref{comp:propagation}; every sufficiently late
fixed-offset law passes its tests.
The same common-word coupling and intrinsic indexed-kernel identities
identify these offsets with the physical annular family, as required
by Proposition~\ref{comp:relative-units}.
Apply Proposition~\ref{comp:relative-units}, then
Proposition~\ref{comp:exact-power-comparison} and rigidity, to any
subsequential family.  The quantile bounds give the stated interval
for $C$.

For arbitrary time units $n_r\to\infty$, put
\[
 k_r=\lfloor\log_4 n_r\rfloor,\qquad m_r=4^{k_r},\qquad
 u_r=n_r/m_r.
\]
Pass to a subsequence with $u_r\to u\in[1,4]$.
Before extracting the $m_r$-time arrays, add the countable family of
fixed $u4^m$ transforms, $m\in\mathbb Z$, of their charts and reference
experiments.  Retain also the actual and reference experiments in
$n_r$ time units, sampling their reference duration and shift densities
in those target units as in Proposition~\ref{loc:kernel-closure}.
Retain the full joint continuum input and its fixed-$u$
transform, including the actual fields, curves, coordinate dilation and
internal reference metrics.  Lemma~\ref{comp:quantile-couplings}, the
exact joint transport in Lemma~\ref{comp:fk-reference-atlas}, and
Proposition~\ref{loc:kernel-closure} permit this enlargement with one
common residual rotation and one transported imaginary-field completion.
The joint background law is fixed along the extraction; no continuity
of its measurable reembedding as a function of $u_r$ is required.
If $Z_r$ is the retained normalized base contour, the contour in $n_r$
units is approximated by $u_r^{-1/2}Z_r(u_r\,\cdot)$.  The contour
modulus identifies its limit with the fixed-$u$ Brownian rescaling.

Keep graph costs in the units $A_{k_r}=B(n_r)$.  In the fixed-$u$
coordinates, let $\phi_{u,r}$ map back to the base coordinates and let
$\eta_r$ be the retained base curve.  The projection at $n_r$-time label
$t$ is
\[
 \phi_{u,r}^{-1}\eta_r(u_rt)
       =\eta_r^{(u)}\bigl((u_r/u)t\bigr),\qquad
 \eta_r^{(u)}(t)=\phi_{u,r}^{-1}\eta_r(ut).
\]
The fixed-$u$ curve law has a tight local uniform modulus, so this
projection differs uniformly on compact time windows from
$\eta_r^{(u)}(t)$ by a quantity tending to zero in probability.
Retain the raw marked path arrays on fixed enlarged windows.  Their
graph paths and costs are unchanged; only their time labels and spatial
projections have been relabeled.  Strict open confinement and port
margins preserve upper realizations, and the reverse marked extraction
retains every bounded-cost confined path.  Thus the limiting passage
collection is the fixed-$u$ transport of the power-of-four extraction.
This step uses geometric margins for certificates, not a bound for
joining prescribed discrete endpoints.

If the latter extraction has multiplier $C$, its transported passage
rule, with the graph costs still in the base units, satisfies
\[
 \widetilde F_U(x,y)
   =C D_h(\phi_u(x),\phi_u(y);\phi_u(U))
   =Cu^pD_{h^{(u)}}(x,y;U).
\]
The multiplier $Cu^p$ is deterministic and remains in a fixed compact
subset of $(0,\infty)$ because $u\in[1,4]$.  Intrinsic scope follows
from Lemma~\ref{comp:intrinsic-lower} and
Equation~\eqref{comp:upper-internal}.
\end{proof}
